\documentclass[11pt,a4paper]{article}
\usepackage[utf8]{inputenc}
\usepackage[T1]{fontenc}
\usepackage[margin=1in]{geometry}
\usepackage{xcolor}
\usepackage{hyperref}
\usepackage{setspace}
\setstretch{1.15}

\hypersetup{colorlinks=true,linkcolor=blue!70!black,citecolor=blue!70!black,urlcolor=blue!70!black}

\title{\textbf{Point-by-Point Response to Technical Check}\\[0.3em]
\large \textit{npj Artificial Intelligence}}
\author{\textbf{Manuscript ID:} 5a3d24b0-5398-4e89-835d-34efbe8bd9e3\\
\textbf{Revision:} v2.3\\
\textbf{Title:} Financial Time-Series Generation: A Systematic Survey of Taxonomy, Evaluation, Datasets, and Open Challenges\\
\textbf{Authors:} Yingxiao Zhang, Jiaxin Duan, Yue Zou, Ke Feng, Jing Xiao\\
\textbf{Assistant Editor:} Dr.~Monika Samra}
\date{September 10, 2026}

\begin{document}

\maketitle

\noindent Dear Dr.~Monika Samra and Editorial Team of \textit{npj Artificial Intelligence},

\vspace{0.8em}
\noindent We thank the editorial team for the further technical check of our manuscript. We have addressed both points raised in the most recent Technical Check letter. As instructed, no changes have been made beyond the requested items. The clean revised manuscript is provided as \texttt{is26.pdf}, and the marked-up manuscript showing every change in red font is provided as \texttt{is26\_marked.pdf} under `Related Files'. Both files compile without errors, with all citations and cross-references resolving correctly and no undefined references; the clean manuscript is 47 pages and the marked-up manuscript is 50 pages.

\vspace{1.5em}
\hrule
\vspace{1.5em}

\subsection*{Point 1: Results and Methods must be structured using only primary subheadings}

\vspace{0.5em}
\noindent\textbf{Editor's Comment:}
\textit{``Please note that Results and Methods sections should be structured using only primary subheadings. We do not permit the use of secondary subheadings. Any bolded or numbered text placed under primary headings, or any other subsection made under primary subheadings are considered as secondary subheadings. Please edit accordingly.''}

\vspace{0.5em}
\noindent\textbf{Response:}
We accept this point and have applied it to the letter of the instruction. We interpreted the rule as having three consequences, and we addressed all three: (i) every heading below the level of the compulsory primary subheadings in Results and Methods was deleted; (ii) every bolded list label and every numbered or bulleted list in those two sections was dismantled and rewritten as continuous running prose, since bolded or numbered text under a primary heading is itself classified as a secondary subheading; and (iii) all textual cross-references to former ``subsections'' were rewritten, so that no reference now points to a heading level that no longer exists.

\vspace{0.5em}
\noindent\textbf{(i) Secondary subheadings removed.} In the previous revision, 38 secondary subheadings remained under the Results section, and a number of headings carried a label but no text of their own. All of them have now been deleted. Five headings that had no content beneath them have been removed outright. At each site where a secondary subheading was deleted, a red bracketed marker of the form \textcolor{red}{[Secondary subheading ``\ldots'' removed.]} is placed in the marked-up manuscript, so that each deletion can be checked individually against this letter. The primary subheadings that remain in Results are exactly the following seven, all of which are the compulsory IRDM subheadings requested by the journal:

\begin{itemize}
  \item Problem Formulation and A Two-Level Taxonomy
  \item Financial Time-Series Extrapolation
  \item Financial Time-Series Imputation
  \item Financial Time-Series Synthesis
  \item Applications
  \item Evaluation Protocols
  \item Data Resources and Benchmarks
\end{itemize}

\noindent In Methods, the four compulsory primary subheadings are retained and no secondary subheading remains: \textit{Literature Retrieval Strategy}, \textit{Inclusion and Exclusion Criteria}, \textit{Data Extraction and Coding Protocol}, and \textit{Review Dimensions and Taxonomy Synthesis}.

\vspace{0.5em}
\noindent\textbf{(ii) All bolded and numbered text converted to prose.} Because bolded labels or numbering placed under a primary heading are treated as secondary subheadings, the lists that previously appeared in Results and Methods have been removed entirely. In total, 14 itemized or enumerated lists were dismantled, covering 49 bolded list labels. Each list item has been rewritten as one or more complete sentences within a continuous paragraph, with the former bold label incorporated as a natural subject of the sentence rather than as a lead-in label. For example, a former entry reading ``\textbf{Temporal Dependence:} Financial data is fundamentally sequential\ldots'' is now rendered as ``Temporal dependence matters because financial data are fundamentally sequential\ldots''. Every sentence created by this rewriting is shown in red in the marked-up manuscript. As a result, no bold text, no bullet marker, and no numbering now appears anywhere under any primary heading in Results or Methods; the sections consist solely of continuous prose paragraphs beneath the primary subheadings.

\vspace{0.5em}
\noindent\textbf{(iii) Cross-references corrected.} Textual references that previously pointed to ``subsections'' (for instance in the Introduction, the Methods section on review dimensions, the Applications section, and the Data Availability declaration) have been rewritten to refer to sections, so that the text no longer refers to a heading level that has been abolished.

\vspace{0.5em}
\noindent\textbf{Discussion unchanged.} The Discussion remains a sequence of paragraphs without subheadings of any level, as required, and its content has not been altered in this revision.

\vspace{1.5em}
\hrule
\vspace{1.5em}

\subsection*{Point 2: Marked-up manuscript showing changes in red font}

\vspace{0.5em}
\noindent\textbf{Editor's Comment:}
\textit{``Marked-up Manuscript: Please include a marked-up manuscript, showing changes to the text in red font. Please do not include tracked changes in comments as we're unable to view them and upload it under 'Related Files' section in pdf or docx. Format.''}

\vspace{0.5em}
\noindent\textbf{Response:}
We have provided a complete marked-up manuscript, \texttt{is26\_marked.pdf}, which has been uploaded under the `Related Files' section in PDF format as requested. The file satisfies each element of the instruction.

\vspace{0.5em}
\noindent\textbf{Changes shown in red font.} Every change made in this revision is rendered directly in red font, so that it is visible in any standard PDF reader without any markup tooling. A total of 111 red-font sites now appear across 28 pages of the manuscript. All newly written prose (the former list content of Results, the former list content of Methods, and the reworded sentences) is printed in red. In addition, red bracketed markers of the form \textcolor{red}{[Secondary subheading ``\ldots'' removed.]} are placed at each of the 47 positions where a secondary subheading was deleted, so that the editor can verify the restructuring point by point and confirm that no secondary heading has been overlooked.

\vspace{0.5em}
\noindent\textbf{No tracked changes and no comments.} In accordance with your instruction that tracked changes in comments cannot be viewed, the marked-up file contains no tracked-change records and no comment or annotation balloons. The red font and red bracketed markers are typeset directly into the document text; nothing about the revision is conveyed through a PDF annotation layer. The file can therefore be read and verified in any viewer.

\vspace{0.5em}
\noindent\textbf{Explanatory summary box.} The marked-up manuscript carries a boxed summary of the revisions, divided into (A) the round-2 structural formatting changes and (B) the round-3 changes to Results and Methods. This box states precisely how many subheadings were removed, how many lists were converted, and which primary subheadings remain, providing the editorial office with a checklist-level overview that can be compared item by item against the red markers in the text.

\vspace{0.5em}
\noindent\textbf{Clean manuscript also provided.} The clean manuscript, \texttt{is26.pdf}, contains no red font and no revision markers; it shows only the final text as it is intended to appear in production. Red bracketed markers are generated by a conditional macro that produces no output in the clean build, so the two files are guaranteed to be otherwise identical in content.

\vspace{1.5em}
\hrule
\vspace{1.5em}

\noindent\textbf{Files accompanying this revision:}

\begin{itemize}
  \item \texttt{is26.pdf} --- clean revised manuscript (no red marks).
  \item \texttt{is26\_marked.pdf} --- marked-up manuscript, all changes in red font, no tracked changes or comments, uploaded under `Related Files'.
  \item \texttt{Response\_to\_Technical\_Check\_Round3.pdf} --- this point-by-point response.
  \item \texttt{Supplementary\_Information\_PRISMA.pdf} --- PRISMA 2020 checklist, with manuscript locations updated to the current primary headings.
  \item \texttt{Figure1.pdf}--\texttt{Figure6.png} --- individual figure files without embedded labels or titles.
\end{itemize}

\vspace{1em}
\noindent We are grateful for the editorial team's patience and guidance, and we hope that the manuscript now satisfies the technical requirements in full.

\vspace{1.6em}
\noindent Sincerely,

\vspace{1em}
\noindent\textbf{Ke Feng, Ph.D.} (Corresponding Author)\\
School of Software and Microelectronics, Peking University, Beijing, China\\
E-mail: \texttt{fengke@pku.edu.cn}\\[0.8em]
\textbf{Jing Xiao, Ph.D.} (Corresponding Author)\\
Chief Scientist, Pingan Group Co., Ltd., Shenzhen, China\\
E-mail: \texttt{xiaojing661@pingan.com.cn}

\end{document}
